%%%
%%%
%%%
%%%
\section{The main classes / namespaces}
Please also read Subsections \ref{subsec_main_principles}, \ref{subsec_glossary}, \ref{subsec_overview_folder_structure}.
The main classes are the following. Parentheses comprise the folder where the .h files for these classes can be found.
\begin{description}
\item[(common/geometry/) swMesh]
\item [(common/) swCell] Attributes and methods for Cell-type objects
\item [(common/) swEdgeBoundary] Attributes and methods for boundary edge-type objects
\item [(common/) swEdgeIntern] Attributes and methods for internal edge-type objects
\item [(common/) swLog] Attributes and methods for log file writing
\item [(common/) swNode] Attributes and methods for node-type objects
\item [(common/) swIO] Contains the attributes and methods for Input-Output purposes.
\item [(solver/models/common) swBoundaryConditionRatingCurves] Methods and attributes for handling rating curve-like boundary conditions
\item [(solver/models/common) swBoundaryConditionTimeSeries] Methods and attributes for handling time series-like boundary conditions
\item [(solver/models/common) swSource] Methods and attributes for handling source terms
\item [(solver/models/) swParent]
\item [(solver/models/) swSwes]
\end{description}
%%%
%%%
%%%
%%%
\section{Variable types, naming conventions}
The main variable types are the following:
\begin{itemize}
\item cell-based variables
\item (boundary) edge-based variables
\item (internal) edge-based variables
\item node-based variables
\end{itemize}
The 
%%%
%%%
%%%
%%%
\section{Solver: code architecture}
%%
%%
The structure is the following:
\begin{verbatim}
main
|- swParent
|- sw<Model>
    |- sw<Model>.h
    |- sw<Model><Process><Type><Sub-type><Method>.cpp
\end{verbatim}
%%%
%%%
%%%
%%%
\section{Solver: simulation sequence}
%%%
%%%
\subsection{main}
The \textbf{main} performs the following tasks:
\begin{itemize}
\item gets the name of the \swd project folder (where the /input and other folders are located)
\item instantiates the \texttt{swIO} class, so as to make I/O methods available + define the variable \texttt{simulation\_folder}. The pointer is \texttt{m\_io}.
\item creates log solver file (log/consoleSolver.log)
\item Read the solver parameters (\texttt{swSolverParam} class). The pointer is \texttt{m\_param\_solver}
\item Creates the Mesh (class \texttt{swMesh}). The pointer is \texttt{m\_mesh}. Builds the mesh, i.e. get the cell/edge/node properties and mesh connectivity info (reading instruction in \texttt{readGeometriyFile*.cpp}).
\item Select the model and create it (in \texttt{swParent} class). The \texttt{solver\_param} pointer is an input argument to the method. 
\item Call \texttt{m\_model->init()}
\item Call \texttt{m\_model->execute()}


\end{itemize}
%%%
%%%
\subsection{\texttt{sw<Model>.init()}}
This method performs a number of model-dependent tasks. Detailed here for Swes.
%%
%%
\begin{itemize}
%
\item Imports the solver parameters required by the model from the \texttt{m\_solver\_param->} attributes.
%
\item Builds up the mesh properties required by the model. This is done by the \texttt{setMeshTables()} method.
%
\item Reads and assigns the initial conditions. These are model-specific. Their number and nature does not depend on the processes that are activated.
%
\item Reads and assigns the boundary conditions. These are model-specific. Their number and nature does not depend on the processes that are activated. 
%
\begin{enumerate}
\item Sets up the following BC pointers: \texttt{bc\_ratingcurve\_list} and \texttt{bc\_timesries\_list}.
\item uses these pointers to read the BC files (using the appropriate m\_io pointer) by calling \texttt{readBoundaryTimeSeriesFile()}.
\item sets up unorderd maps so as to associate the BC codes (that are negative) with the corresponding time series. At the end of \texttt{intBoundaryCondition}, for each boundary edge $i$, the pointer \texttt{edge\_bnd\_bc\_pointer[i]} contains a pointer called \texttt{bc} of type \texttt{swBoundaryCondition} that provides access to the BC time series and type.
\end{enumerate}
%
\item Reads and assigns the various model parameters. These are specific to the local \texttt{sw<Model>} class. The number and type of parameters to be read my depend on some solver parameters. For instance, if a specific friction model is selected (\textbf{model\_friction $\neq$ "none"}), the method that reads the friction parameter from the friction map file and assigns it to the appropriate parameter must be called.
%
\item (optional) Compute the edge factors that are used by the \texttt{computeBalance()} method for the sake of faster CPU
%
\item resize the fluxes and Riemann state vectors
%
\item Writes (text) log fils in the \textbf{/log} folder to serve as a tracebak for solver execution.
\end{itemize}
%%%
%%%
\subsection{\texttt{sw<Model>.execute()}}
%%%
%%%

All models are different by definition. However, all transient shallow water models involve the following compulsory steps:
\begin{itemize}
%
%
\item Prior to time loop: initialize simulation time, etc.
%
%
\item Time loop: 
\begin{enumerate}
%
\item \texttt{computeBC()}:interpolates the boundary conditions from the tables initialized at the \texttt{sw<Model>.init()} stage.
%
\item \texttt{hypCompute()} : computes the hyperbolic terms for the hydrodynamic modeln. Computing them using explicit finite volumes requires the following steps:
\begin{enumerate}
%
\item \texttt{computeBC()}:interpolates the boundary conditions from the tables initialized at the \texttt{sw<Model>.execute()} stage.
%
\item \texttt{computeLRStates()}: Compute the left and right states of the Riemann problems: This step is method-dependent (the computation is not the same if a Godunov or a MUSCL scheme is used)
%
\item \texttt{computeTimeStep()}. The storage period and CFL constraints require that $\Delta t$ be computed at each time level.
%
\item \texttt{computeFluxes()}. The fluxes across boundary and internal interfaces are computed by solving the Riemann problems.
%
\item \texttt{computeBalance()}. The balance is carried out for mass, momentum and any other conserved variable, using the previously computed fluxes as a basis.
%
\item \texttt{DivCorr()}. Apply divergence correction to avoid cell over-drying.
%
\end{enumerate}
%
%
\item \texttt{parabCompute()}. Compute the contribution of parabolic terms (if applicable)
%
%
\item Source terms. The list of applicable source terms is model-dependent. For instance, in the SWES equations, the possible source terms are friction, wind-induced stress, precipitation/infiltration. They are applied sequentlially, following a first-order time splitting.

%

\end{enumerate}
%
%
%
\item After time loop: close files, finalize CPU time computation
\end{itemize}
%%%
%%%
%%%
\section{Where can I find / What are these variables}
%%%
%%%
\subsection{pointers}
\begin{description}
\item [bc\_ratingcurve\_list]: \texttt{swBoundaryConditionRating}. Pointer defined at the \texttt{initBoundaryConditions()} stage (belongs to \texttt{swParent}).
\item [bc\_timeseries\_list]: \texttt{swBoundaryConditionTimeSeries}. Pointer defined at the \texttt{initBoundaryConditions()} stage (belongs to \texttt{swParent}).
\item [m\_io] : \texttt{swIO} class. Pointer defined at the class construction stage. The class is instantiated in the main.cpp of swSolver. \\
How to get it from a given model : \texttt{m\_solver\_param->m\_io->}
\item [m\_mesh] : \texttt{swMesh} class. Pointer defined at the class construction stage. The class is instantiated in the main.cpp of swSolver.
\item [m\_model]: \texttt{sw<Model>} class. Pointer defined by \texttt{swParent::create()}
\item [m\_solver\_param] : \texttt{swSolverParam} class. Pointer defined at the class construction stage. The class is instantiated in the main.cpp of swSolver. 
\end{description}
%%%
%%%
\subsection{File names}
%%%
\begin{description}
\item [Reserved file names]: these are provided in the \texttt{src/common/swFile.h}. They are relative to the root of the project folder.
\item [swIO::getAbsolutePath]: this method takes a relative path (relative to the simulation folder) and as an input and returns the absolute path.\\
Example : if the user runs 
\begin{verbatim}
sw2dSolver test/test1
\end{verbatim}
then swIO::getAbsolutePath("input/geometry.sw2d") will return "test/test1/input/geometry.sw2d"
\end{description}
%%%
%%%
